\subsection{SIMPLE MODULES WITH A HIGHEST WEIGHT}

Recall that fixing B defines an order relation on \(h_{Q}^{*}\) (Chap. VI, §1, no. 6). The elements of \(h_{Q}^{*}\) that are \(\geq 0\) are the linear combinations of elements of B with rational coefficients \(\geq 0\) .

More generally, we shall consider the following order relation between elements \(\lambda, \mu \in \mathfrak{h}^*\) :

\(\lambda - \mu\) is a linear combination of elements of B with rational coefficients \(\geq 0\) .

Lemma 2. Let V be a simple g-module, \(\omega\) a weight of V. The following conditions are equivalent:

\begin{enumerate}
\def\labelenumi{(\roman{enumi})}
\item
  every weight of V is of the form \(\omega -\mu\) where \(\mu\) is a radical weight \(\geq 0\) ;
\item
  \(\omega\) is the highest weight of V;
\item
  for all \(\alpha \in \mathrm{B}\) , \(\omega + \alpha\) is not a weight of V;
\item
  there exists a primitive element of weight \(\omega\) .
\item
  \(\Longrightarrow\) (ii) \(\Longrightarrow\) (iii): This is clear.
\item
  \(\Longrightarrow\) (iv): Assume that condition (iii) is satisfied. For all \(h\in \mathfrak{h}\) ,
\end{enumerate}

\(\operatorname {Ker}(h_{\mathrm{V}} - \omega (h))\)

is non-zero, contained in \(V^{\omega}\) , and stable under \(h_{V}\) . By induction on dim h, we see that there exists a non-zero v in \(V^{\omega}\) such that \(h v \subset kv\) . Condition (iii) implies that \(n_{+}v = 0\) , so v is primitive.

\begin{enumerate}
\def\labelenumi{(\roman{enumi})}
\setcounter{enumi}{3}
\tightlist
\item
  \(\Longrightarrow\) (i): Let \(v\) be a primitive element of weight \(\omega\) . Since V is simple, V is generated by \(v\) as a \(\mathfrak{g}\) -module. Assertion (i) now follows from Prop. 1.
\end{enumerate}

Q.E.D.

Thus, for any simple g-module, the existence of a primitive element is equivalent to that of a highest weight, or to that of a maximal weight.

There exist simple \(\mathfrak{sl}(2,\mathbf{C})\) -modules V that have no weights for any Cartan subalgebra \(\mathfrak{h}\) of \(\mathfrak{sl}(2,\mathbf{C})\) (§1, Exerc. 14 \(f\) )). These modules are of infinite dimension over \(\mathbf{C}\) (§1, no. 3, Th. 1).

PROPOSITION 3. Let V be a simple \(\mathfrak{g}\) -module with a highest weight \(\omega\) .

\begin{enumerate}
\def\labelenumi{(\roman{enumi})}
\item
  The primitive elements of V are the non-zero elements of \(V^{\omega}\) .
\item
  V is semi-simple as an \(\mathfrak{h}\) -module.
\item
  We have \(V = \bigoplus_{\lambda \in h^{*}} V^{\lambda}\) . For all \(\lambda \in h^{*}\) , \(V^{\lambda}\) is finite dimensional. We have \(\dim V^{\omega} = 1\) .
\item
  The \(\mathfrak{g}\) -module V is absolutely simple.
\end{enumerate}

Assertions (i), (ii) and (iii) follow from Prop. 1 and Lemma 2. Assertion (iv) follows from Prop. 1 (iii) and Algebra, Chap. VIII, §7, no. 3.

COROLLARY. If \(\mathrm{V}\) is finite dimensional, the canonical homomorphism \(\mathrm{U}(\mathfrak{g}) \to \operatorname{End}(\mathrm{V})\) is surjective.

This follows from (iv), cf.~Algebra, Chap. VIII, §3, no. 3.

PROPOSITION 4. Let V be a simple g-module with a highest weight \(\omega\) , X a semi-simple g-module, and \(X'\) the isotypical component of type V in X. Then \(X'\) is the submodule of X generated by \(X_{\pi}^{\omega}\) . Its length is equal to the dimension of \(X_{\pi}^{\omega}\) .

Let \(X''\) be the submodule of X generated by \(X_{\pi}^{\omega}\) . It is clear that every submodule of X isomorphic to V is contained in \(X''\) . Hence \(X' \subset X''\) . On the other hand, let \(\varPhi\) be the set of homomorphisms from the g-module V to the g-module X. The length of \(X'\) is \(\dim_k \varPhi\) (Algebra, Chap. VIII, §4, no. 4), that is \(\dim_k X_{\pi}^{\omega}\) (Prop. 2).

